\section{Simulation}
\label{sec:simulations}

Numerical simulations are conducted to verify the effectiveness of the proposed RBF neural network-based adaptive control strategy. 
Two nonlinear systems with different uncertainty characteristics are considered. 
The first example investigates the approximation capability of the RBF compensator under non-smooth disturbances, while the second example evaluates the proposed controller for a multi-joint robotic manipulator with dynamic coupling and lumped uncertainties.

\begin{table}[!t]
	\caption{Definitions of Control Modes and Adaptive Laws}
	\label{tab:control_modes}
	\centering
	\begin{tabular}{c c l}
		\hline
		Parameter & Value & Description \\
		\hline
		$S_1$ & $1$ & First adaptive law \\
		$S_1$ & $2$ & Second adaptive law with $\sigma$-modification \\
		\hline
		$S$ & $1$ & Nominal model-based controller without compensation \\
		$S$ & $2$ & Controller with exact uncertainty compensation \\
		$S$ & $3$ & Controller with RBF neural network compensation \\
		\hline
	\end{tabular}
\end{table}

%======================================================
%======================================================
\subsection{Example 1: Nonlinear Servo System with Non-Smooth Disturbance}

Consider a single-degree-of-freedom nonlinear servo system described by
\begin{equation}
	M\ddot{q}=\tau+d(\dot{q}),
\end{equation}
where $M=10$ and the disturbance term is given by
\begin{equation}
	d(\dot{q})=-15\dot{q}-30\operatorname{sgn}(\dot{q}).
\end{equation}

Before presenting the simulation results, the overall simulation
configuration is introduced. The Simulink implementation of the proposed
control framework is illustrated in Fig.~\ref{fig:sim_structure}.
The simulation platform consists of the nonlinear servo system, the
controller module, the reference trajectory generator, and the uncertainty
compensation module.

The controller receives the desired trajectory and the measured system
states, and generates the control input $\tau$. Depending on the selected
control mode, the uncertainty compensation term is either disabled
($S=1$), obtained from the known disturbance model ($S=2$), or generated
by the RBF neural network ($S=3$).

\begin{figure}[!t]
	\centering
	\includegraphics[width=0.48\textwidth]
	{result_pdf/chap6_1sim.pdf}
	\caption{
		Simulink implementation of the nonlinear servo system and the proposed
		control framework.
	}
	\label{fig:sim_structure}
\end{figure}

The disturbance consists of viscous friction and Coulomb friction components, 
which introduces a non-smooth uncertainty into the closed-loop system.

The desired trajectory is selected as
\begin{equation}
	q_d(t)=\sin(t)\ \mathrm{rad},
\end{equation}
with the initial condition
\[
[q(0),\dot q(0)]^T=[0.6,0]^T .
\]


The controller parameters are selected as
\[
Q=50I_2,\qquad \gamma=200,\qquad \sigma=0.001 .
\]

The RBF neural network employs Gaussian basis functions with different
numbers of hidden nodes. Two configurations with $m=5$ and $m=19$ are
considered.

For the case of $m=5$, the centers are selected as

\[
c_i=
[\alpha_i,\alpha_i]^T,
\]

where

\[
\alpha_i\in\{-1,-0.5,0,0.5,1\}.
\]

For the case of $m=19$, a denser distribution of centers is adopted:

\[
\alpha_i\in
\{-1,-8/9,\ldots,0,\ldots,8/9,1\}.
\]

The Gaussian basis functions are defined as

\[
h_i(\xi)=
\exp
\left(
-\frac{\|\xi-c_i\|^2}{2b^2}
\right),
\]

where the width parameter is selected as $b=5$ for all basis functions.
The initial neural network weights are set as
$\hat W(0)=0$.


To evaluate the effectiveness of the proposed adaptive compensation strategy,
three different compensation modes are considered:

\begin{itemize}
	\item $S=1$: nominal model-based controller without uncertainty compensation;
	\item $S=2$: exact compensation controller assuming that the uncertainty 
	$d(\dot q)$ is completely known;
	\item $S=3$: RBF neural network compensation controller, where the unknown
	uncertainty is approximated online.
\end{itemize}

For $S=3$, two adaptive laws are investigated:
$S_1=1$ and $S_1=2$.


%======================================================
\subsubsection{Nominal Model-Based Controller ($S=1$)}

In this case, only the nominal model is used in the controller design.
The disturbance term is not compensated, and the adaptive law selection
parameter $S_1$ has no influence because the RBF network is not involved.

The tracking response is shown in Fig.~\ref{fig:S1_track}, while the
uncertainty response is presented in Fig.~\ref{fig:S1_rbf}.

\begin{figure}[!t]
	\centering
	\includegraphics[width=0.48\textwidth]
	{result_pdf/chap1_S=1_displacementandvelocity.pdf}
	\caption{
		Tracking performance of the nominal controller ($S=1$).
	}
	\label{fig:S1_track}
\end{figure}


\begin{figure}[!t]
	\centering
	\includegraphics[width=0.48\textwidth]
	{result_pdf/chap1_S=1_RBFfx.pdf}
	\caption{
		Unknown disturbance and estimation result under the nominal controller
		($S=1$).
	}
	\label{fig:S1_rbf}
\end{figure}


As observed in Fig.~\ref{fig:S1_track}, the nominal controller cannot
completely reject the non-smooth friction disturbance, resulting in a
non-zero tracking error.

Since no compensation mechanism is activated, the estimated uncertainty
remains zero, as shown in Fig.~\ref{fig:S1_rbf}.


%======================================================
\subsubsection{Exact Compensation Controller ($S=2$)}

For $S=2$, the uncertainty term is assumed to be known exactly.
Therefore, the compensation input directly cancels the disturbance:

\[
\tau_c=-d(\dot q).
\]

The adaptive law parameter $S_1$ does not affect the control performance
because the RBF neural network is not used.


\begin{figure}[!t]
	\centering
	\includegraphics[width=0.48\textwidth]
	{result_pdf/chap1_S=2_displacementandvelocity.pdf}
	\caption{
		Tracking performance with exact uncertainty compensation ($S=2$).
	}
	\label{fig:S2_track}
\end{figure}


\begin{figure}[!t]
	\centering
	\includegraphics[width=0.48\textwidth]
	{result_pdf/chap1_S=2_RBFfx.pdf}
	\caption{
		Exact disturbance compensation result ($S=2$).
	}
	\label{fig:S2_rbf}
\end{figure}


As shown in Fig.~\ref{fig:S2_track}, the tracking error converges rapidly
because the disturbance is completely compensated.

Different from the RBF compensation case, the uncertainty estimation is
not generated by the neural network. Instead, the exact disturbance model
is directly used for compensation.

%======================================================
\subsubsection{RBF Neural Network Compensation ($S=3$)}

When $S=3$, the explicit expression of the nonlinear disturbance is assumed
to be unknown to the controller. Therefore, an RBF neural network is
introduced to approximate the lumped uncertainty:

\begin{equation}
	\hat{f}(x)=\hat{W}^{T}h(\xi),
\end{equation}

where $\hat{W}$ denotes the adaptive weight vector and $h(\xi)$ represents
the Gaussian basis function vector.

Unlike the previous two cases, the disturbance compensation is generated
by the neural network rather than directly using the known disturbance
model. Therefore, the adaptive law plays an important role in determining
the convergence behavior of the estimated uncertainty.


%------------------------------------------------------
\paragraph{First adaptive law ($S=3,S_1=1$)}

Firstly, the RBF controller with the first adaptive law is investigated.
The tracking performance and uncertainty approximation results are shown
in Fig.~\ref{fig:S31_track} and Fig.~\ref{fig:S31_rbf}, respectively.


\begin{figure}[!t]
	\centering
	\includegraphics[width=0.48\textwidth]
	{result_pdf/chap1_S=3_S1=1_displacementandvelocity.pdf}
	\caption{
		Tracking performance of the RBF adaptive controller with the first
		adaptive law ($S=3,S_1=1$).
	}
	\label{fig:S31_track}
\end{figure}


\begin{figure}[!t]
	\centering
	\includegraphics[width=0.48\textwidth]
	{result_pdf/chap1_S=3_S1=1_RBFfx.pdf}
	\caption{
		Approximation of the nonlinear disturbance using the RBF neural network
		with the first adaptive law ($S=3,S_1=1$).
	}
	\label{fig:S31_rbf}
\end{figure}


As shown in Fig.~\ref{fig:S31_track}, the RBF compensation effectively
reduces the tracking error compared with the nominal controller.
The approximation result in Fig.~\ref{fig:S31_rbf} demonstrates that the
neural network is able to learn the unknown disturbance online.


%------------------------------------------------------
\paragraph{Second adaptive law with $m=5$ hidden nodes
	($S=3,S_1=2$)}

Next, the second adaptive law is evaluated with an RBF neural network
containing five hidden nodes.

The corresponding tracking response and disturbance approximation result
are shown in Fig.~\ref{fig:S32_m5_track} and Fig.~\ref{fig:S32_m5_rbf}.


\begin{figure}[!t]
	\centering
	\includegraphics[width=0.48\textwidth]
	{result_pdf/chap1_S=3_S1=2_displacementandvelocity_m=5.pdf}
	\caption{
		Tracking response using the second adaptive law with five hidden nodes
		($S=3,S_1=2,m=5$).
	}
	\label{fig:S32_m5_track}
\end{figure}


\begin{figure}[!t]
	\centering
	\includegraphics[width=0.48\textwidth]
	{result_pdf/chap1_S=3_S1=2_RBFfx_m=5.pdf}
	\caption{
		RBF approximation result with five hidden nodes
		($S=3,S_1=2,m=5$).
	}
	\label{fig:S32_m5_rbf}
\end{figure}


With a limited number of hidden nodes, the RBF network provides a coarse
approximation of the nonlinear disturbance. The residual approximation
error leads to a relatively larger tracking deviation.


%------------------------------------------------------
\paragraph{Second adaptive law with $m=19$ hidden nodes
	($S=3,S_1=2$)}

To further investigate the influence of the approximation capability,
the number of hidden nodes is increased to $m=19$.

The corresponding results are presented in
Fig.~\ref{fig:S32_m19_track} and Fig.~\ref{fig:S32_m19_rbf}.


\begin{figure}[!t]
	\centering
	\includegraphics[width=0.48\textwidth]
	{result_pdf/chap1_S=3_S1=2_displacementandvelocity_m=19.pdf}
	\caption{
		Tracking response using the second adaptive law with nineteen hidden
		nodes ($S=3,S_1=2,m=19$).
	}
	\label{fig:S32_m19_track}
\end{figure}


\begin{figure}[!t]
	\centering
	\includegraphics[width=0.48\textwidth]
	{result_pdf/chap1_S=3_S1=2_RBFfx_m=19.pdf}
	\caption{
		RBF approximation result with nineteen hidden nodes
		($S=3,S_1=2,m=19$).
	}
	\label{fig:S32_m19_rbf}
\end{figure}


Compared with the case of $m=5$, the denser distribution of Gaussian basis
functions improves the approximation capability of the RBF neural network.
Consequently, the estimated uncertainty is closer to the actual disturbance,
which leads to improved tracking performance.


%------------------------------------------------------
\paragraph{Summary}

The simulation results verify the effectiveness of the RBF-based adaptive
compensation strategy.

When the uncertainty model is unavailable, the RBF neural network can
successfully approximate the unknown nonlinear disturbance and reduce the
tracking error.

Furthermore, increasing the hidden-layer size from $m=5$ to $m=19$
enhances the approximation capability of the neural network.
The comparison also shows that the adaptive law selection affects the
weight update behavior and uncertainty estimation performance.


%======================================================
\subsubsection{Discussion}

The simulation results verify the effectiveness of different compensation
strategies.

The nominal controller ($S=1$) suffers from residual tracking errors because
the disturbance is completely uncompensated.

The exact compensation controller ($S=2$) achieves the best tracking
performance since the disturbance is assumed to be completely known.

For the practical case where the uncertainty model is unavailable,
the RBF-based adaptive controller ($S=3$) successfully approximates the
unknown disturbance and achieves improved tracking performance.

Moreover, the comparison between $S_1=1$ and $S_1=2$ demonstrates that the
second adaptive law provides improved weight regulation while preserving
tracking accuracy.

%======================================================
\subsection{Example 2: Two-Link Robotic Manipulator with Coupled Uncertainties}

A two-link planar robotic manipulator is considered to further evaluate
the proposed adaptive control strategy in the presence of nonlinear
dynamic coupling, parameter uncertainties, and external disturbances.
The simulation configuration implemented in Simulink is shown in
Fig.~\ref{fig:ex2_sim}.


\begin{figure}[!t]
	\centering
	\includegraphics[width=0.48\textwidth]
	{result_pdf/chap6_2sim.pdf}
	\caption{
		Simulink implementation of the two-link robotic manipulator control
		system.
	}
	\label{fig:ex2_sim}
\end{figure}


The nominal dynamics of the robotic manipulator are described as

\begin{equation}
	M(q)\ddot q+C(q,\dot q)\dot q+G(q)=\tau+d ,
	\label{eq:robot_dynamics}
\end{equation}

where the inertia matrix, Coriolis matrix, and gravity vector are given by


\begin{align}
	M(q)
	&=
	\begin{bmatrix}
		v+q_{01}+2q_{02}\cos(q_2)
		&
		q_{01}+q_{02}\cos(q_2)
		\\[2mm]
		q_{01}+q_{02}\cos(q_2)
		&
		q_{01}
	\end{bmatrix},
	\\[3mm]
	C(q,\dot q)
	&=
	\begin{bmatrix}
		-q_{02}\dot q_2\sin(q_2)
		&
		-q_{02}(\dot q_1+\dot q_2)\sin(q_2)
		\\[2mm]
		q_{02}\dot q_1\sin(q_2)
		&
		0
	\end{bmatrix},
	\\[3mm]
	G(q)
	&=
	\begin{bmatrix}
		15g\cos(q_1)
		+
		8.75g\cos(q_1+q_2)
		\\[2mm]
		8.75g\cos(q_1+q_2)
	\end{bmatrix}.
\end{align}


The physical parameters are selected as

\[
v=13.33,\qquad
q_{01}=8.98,\qquad
q_{02}=8.75,
\qquad
g=9.8~\mathrm{m/s^2}.
\]


The external disturbance is selected as an error-dependent nonlinear
uncertainty:

\begin{equation}
	d(t,e,\dot e)
	=
	d_1+d_2\|e\|+d_3\|\dot e\|,
\end{equation}

where

\[
d_1=[2,2]^T,\qquad
d_2=[3,3]^T,\qquad
d_3=[6,6]^T .
\]


To evaluate the robustness of the proposed controller against modeling
uncertainties, a $20\%$ parameter mismatch is introduced into the nominal
dynamic model:

\[
\Delta M=0.2M_0,\qquad
\Delta C=0.2C_0,\qquad
\Delta G=0.2G_0 .
\]


The desired trajectories of the two joints are selected as

\begin{align}
	q_{1d}(t)
	&=
	1+0.2\sin(0.5\pi t),
	\\
	q_{2d}(t)
	&=
	1-0.2\cos(0.5\pi t).
\end{align}


The initial condition is selected as

\[
X(0)
=
[0.6,0.3,0.5,0.5]^T .
\]


The control parameters are chosen as

\[
Q=50I_4,\qquad
\gamma=20,\qquad
\sigma=0.001 .
\]


In this example, the RBF neural network compensation controller
with the second adaptive law ($S=3,S_1=2$) is adopted.
The Gaussian basis functions are initialized with

\[
\hat W(0)=0.1I .
\]

The RBF network approximates the unknown lumped uncertainty as

\[
\hat f(x)=\hat W^Th(x).
\]


The simulation results are presented in the following figures.
The joint angle tracking responses are shown in
Fig.~\ref{fig:ex2_angle}, while the corresponding angular velocity
tracking performance is presented in Fig.~\ref{fig:ex2_velocity}.


\begin{figure}[!t]
	\centering
	\includegraphics[width=0.48\textwidth]
	{result_pdf/chap2_S=3_S1=2 Angle tracking of joint 1 and joint 2.pdf}
	\caption{
		Angle tracking responses of joint 1 and joint 2 under the proposed
		RBF adaptive controller ($S=3,S_1=2$).
	}
	\label{fig:ex2_angle}
\end{figure}


\begin{figure}[!t]
	\centering
	\includegraphics[width=0.48\textwidth]
	{result_pdf/chap2_S=3_S1=2 Angular velocity tracking of joint 1 and joint 2.pdf}
	\caption{
		Angular velocity tracking responses of joint 1 and joint 2.
	}
	\label{fig:ex2_velocity}
\end{figure}


The generated control inputs of the two joints are illustrated in
Fig.~\ref{fig:ex2_control}.


\begin{figure}[!t]
	\centering
	\includegraphics[width=0.48\textwidth]
	{result_pdf/chap2_S=3_S1=2 Control input of joint 1 and joint 2.pdf}
	\caption{
		Control inputs of joint 1 and joint 2 generated by the adaptive
		controller.
	}
	\label{fig:ex2_control}
\end{figure}


Finally, the actual lumped uncertainty and its RBF approximation are
shown in Fig.~\ref{fig:ex2_approx}.


\begin{figure}[!t]
	\centering
	\includegraphics[width=0.48\textwidth]
	{result_pdf/chap2_S=3_S1=2 f(x) of joint 1 and joint 2 and its approximation.pdf}
	\caption{
		Actual lumped uncertainty and its RBF neural network approximation.
	}
	\label{fig:ex2_approx}
\end{figure}


The simulation results demonstrate that the proposed adaptive RBF
compensation controller achieves accurate trajectory tracking despite the
existence of nonlinear coupling, parameter uncertainties, and external
disturbances.

As shown in Fig.~\ref{fig:ex2_angle} and Fig.~\ref{fig:ex2_velocity},
both joint positions and angular velocities converge to their desired
trajectories with small tracking errors.

Furthermore, Fig.~\ref{fig:ex2_approx} verifies that the RBF neural
network can effectively approximate the unknown lumped uncertainty online.
The generated control inputs remain bounded and smooth, as illustrated in
Fig.~\ref{fig:ex2_control}, demonstrating the robustness of the proposed
adaptive compensation strategy.

%=================================================================================

